\documentclass[../main-en.tex]{subfiles}
\begin{document}

\section{Module View}
\label{sec:module-view}

% Dense figure packing, scoped to the four UML sections (restored at the end
% of the mapping section, so the rest of the document is unaffected).
\setcounter{topnumber}{3}
\setcounter{totalnumber}{5}
\renewcommand{\topfraction}{0.9}
\renewcommand{\floatpagefraction}{0.85}

The module view shows how the code is split into modules and which module
depends on which. It is where modifiability is mostly decided.

\subsection{Decomposition}

The hub is implemented in C\#. We decomposed it into six modules, each a
UML package that maps to one C\# namespace under the root namespace
\texttt{SensorHub}. Application holds the entry point and orchestrates the
whole hub, Ingestion provides sensor readings from different sources,
Processing transforms readings through a configurable chain, Output delivers
results to the console dashboard and the log file, Configuration loads the
config file and builds the pipeline from it, and Core holds the data types
shared by all the others. \Cref{fig:module-view} shows the decomposition and
the content of each module.

\begin{figure}[htbp]
  \centering
  \wideimage{uml/module-view/Decom.png}
  \caption{Decomposition view: the six modules and their content.}
  \label{fig:module-view}
\end{figure}

The decomposition follows the data flow, with Application, Configuration
and Core around the three domain modules.

\subsection{Uses view}

A module A \emph{uses} a module B when A needs a correct B to complete its
own work~\cite{bass_software_2021}. \Cref{fig:uses-view} shows the uses
relations of the hub. Application uses everything below it. Configuration
uses Ingestion, Processing and Output because its factories create their
objects, plus Core to turn the sensor type named in the config file into a
\texttt{SensorType} value. Ingestion, Processing and Output use nothing but
Core.

\begin{figure}[htbp]
  \centering
  \wideimage{uml/module-view/Uses.png}
  \caption{Uses view: which module needs which to work correctly.}
  \label{fig:uses-view}
\end{figure}

Three rules follow from the diagram:

\begin{itemize}
  \item Ingestion, Processing and Output never use each other. They only
  meet through Application.
  \item Core uses no other module. A change anywhere else cannot break it.
  \item No dependency cycles. The dependencies go from the top (Application,
  Configuration) down to Core.
\end{itemize}

Each rule is a modifiability tactic from the
course~\cite{wiklund_ch8_nodate}: \emph{restrict dependencies} (a change to
a source cannot propagate to a transformation or a sink), \emph{increase
semantic coherence} (Core holds only plain data, so the most used module is
also the most stable), \emph{encapsulate} (\texttt{Hub} works only through
interfaces) and \emph{use an intermediary} (the factories in Configuration
are the only classes that name concrete implementations, as in lab 2 of the
course).

\subsection{Class view}

The class view shows the classes and interfaces inside each module, in
standard UML notation~\cite{freecodecamp_uml_2021}, one diagram per module
(classes of other modules appear as references). Two conventions are enough
to read them: a directed association is a field (the end reads visibility,
role name, multiplicity, for example \texttt{-\_sources 1..*}), and
aggregation (hollow diamond) means the parts can exist on their own while
composition (filled diamond) means they live and die with the whole.

\Cref{fig:class-application,fig:class-core} show Application and Core. The
decision in Application is \emph{encapsulate}: \texttt{Hub} holds one or
more sources, one chain and one or more sinks, all typed by interfaces, and
runs the loop in \texttt{Run()}. The implementations can change without
touching it, and any of them can be replaced by a stub in
tests~\cite{wiklund_ch12_nodate}. Core has no behavior and no dependency:
\texttt{SensorReading} is one measurement, and it carries its own
\texttt{Alert} objects, so every transformation keeps the same signature and
the sinks can show the alerts together with the reading.

\begin{figure}[htbp]
  \centering
  \wideimage{uml/module-view/ApplicationClassView.png}
  \caption{Class view of Application: \texttt{Hub} and the three interfaces
  it depends on.}
  \label{fig:class-application}
\end{figure}

\begin{figure}[htbp]
  \centering
  \includegraphics[width=0.55\textwidth]{uml/module-view/CoreClassView.png}
  \caption{Class view of Core: the shared data types.}
  \label{fig:class-core}
\end{figure}

\Cref{fig:class-ingestion,fig:class-output} show the two ends of the
pipeline. Ingestion offers two sources behind \texttt{ISensorSource}, a CSV
file reader and a random generator, and Output offers two sinks behind
\texttt{IOutputSink}, a console dashboard and a log file. Both the dashboard
and the smoother below keep state per \texttt{(SensorId, SensorType)}: one
sensor id such as \texttt{room-a} reports temperature, CO$_2$ and humidity,
so the latest value must be tracked per type. \texttt{Close()} lets the log
file be flushed when the hub stops.

\begin{figure}[htbp]
  \centering
  \wideimage{uml/module-view/IngestionClassView.png}
  \caption{Class view of Ingestion: the two sources behind
  \texttt{ISensorSource}.}
  \label{fig:class-ingestion}
\end{figure}

\begin{figure}[htbp]
  \centering
  \wideimage{uml/module-view/OutputClassView.png}
  \caption{Class view of Output: the two sinks behind \texttt{IOutputSink}.}
  \label{fig:class-output}
\end{figure}

\Cref{fig:class-processing} shows Processing. Each step implements
\texttt{ITransformation}. \texttt{UnitConverter} applies a linear conversion
defined in the config file, \texttt{MovingAverageSmoother} averages the last
values, \texttt{ThresholdAlerter} adds warning or critical alerts, one class
per rule (\emph{increase semantic coherence}). The
\texttt{TransformationChain} implements
\texttt{ITransformation} itself and holds its steps by aggregation: a step
can exist without a chain, and a step belongs to at most one chain. This is
the Composite pattern, so the chain can be used anywhere a single step is
expected.

\begin{figure}[htbp]
  \centering
  \wideimage{uml/module-view/ProcessingClassView.png}
  \caption{Class view of Processing: the three steps and the chain that runs
  them.}
  \label{fig:class-processing}
\end{figure}

\Cref{fig:class-configuration} shows Configuration. \texttt{ConfigLoader}
parses the JSON file into a \texttt{HubConfig} made of
\texttt{ComponentConfig} entries, one generic entry (a type name plus
key/value settings) used for every source, transformation and sink, so a new
kind of component needs no new config class. The three factories map the
type name to the concrete class, reading typed settings through the
internal helper \texttt{ConfigSettings}. They are the only classes that name a
concrete source, step or sink, and because the chain and its parameters come
from the file, behavior changes with \texttt{config.json}, not with a
rebuild (\emph{defer binding})~\cite{wiklund_ch8_nodate}.

\landscapefigurenopage{uml/module-view/ConfigClassView.png}{Class view of Configuration: the loader, the config types,
  the three factories and their helper.}{fig:class-configuration}

\Cref{lst:config} shows the \texttt{config.json} delivered with the code,
which is the configuration of the component-and-connector view: two sources,
a chain of four steps (conversion, smoothing, then one threshold alerter for
temperature and one for CO$_2$), two sinks, and the pause after each reading.
Each entry is one \texttt{ComponentConfig}: the \texttt{type} selects the
class in the factory, and the \texttt{settings} are passed to its
constructor. The temperature thresholds are in degrees Fahrenheit, because
the alerter runs after the conversion.

\begin{listing}[htbp]
\begin{minted}[fontsize=\footnotesize, bgcolor=LightGray, breaklines]{json}
{
  "delayMs": 200,
  "sources": [
    { "type": "csv", "settings": { "path": "data/room-a.csv" } },
    { "type": "random", "settings": {
        "sensorId": "sim-01", "sensorType": "Temperature", "unit": "C",
        "minValue": "18", "maxValue": "30", "count": "2160" } }
  ],
  "transformations": [
    { "type": "unitConverter", "settings": { "sensorType": "Temperature",
        "fromUnit": "C", "toUnit": "F", "factor": "1.8", "offset": "32" } },
    { "type": "movingAverage", "settings": { "windowSize": "3" } },
    { "type": "threshold", "settings": { "sensorType": "Temperature",
        "warningThreshold": "80.6", "criticalThreshold": "84.2" } },
    { "type": "threshold", "settings": { "sensorType": "CO2",
        "warningThreshold": "1000", "criticalThreshold": "1200" } }
  ],
  "sinks": [
    { "type": "console" },
    { "type": "logFile", "settings": { "path": "hub.log" } }
  ]
}
\end{minted}
\caption{The \texttt{config.json} delivered with the code: the sources, the
transformation chain in order (with the Celsius to Fahrenheit conversion),
the sinks and the delay.}
\label{lst:config}
\end{listing}

\end{document}
